% [VERIFY] intro -- problem.md (O1-O5, G1-G2, Key Insight), Book.md overview, top claims C01-C03,C08.
% MODALITY GUARD: degradation->optimization is "argue/unlikely/do not formally prove" (C01 interp, O3),
% NOT "ruled out". Keep hedged.
\section{Introduction}
\label{sec:intro}

Deep convolutional neural networks have driven a series of breakthroughs in
image classification~\cite{alexnet2012}. Depth is of crucial importance: leading
results on the challenging ImageNet dataset exploit very deep models, and the
features they learn span low/mid/high-level abstractions that many other visual
recognition tasks have benefited from~\cite{zeiler2014}. Driven by this
significance, a natural question arises: \emph{is learning better networks as
easy as stacking more layers?}

An obstacle to answering this question is the notorious problem of vanishing /
exploding gradients~\cite{bengio1994}, which hampers convergence from the
beginning. This problem, however, has been largely addressed by normalized
initialization~\cite{glorot2010,lecun1989,saxe2013} and intermediate
normalization layers, which enable networks with tens of layers to start
converging for stochastic gradient descent (SGD) with backpropagation. When
deeper networks are able to start converging, a \degrad{} has been
exposed~\cite{constrainedtime2015}: with the network depth increasing, accuracy
gets saturated and then degrades rapidly. Unexpectedly, such degradation is
\emph{not} caused by overfitting, and adding more layers to a suitably deep
model leads to \emph{higher training error}. On CIFAR-10, a 56-layer
\plainnet{} reaches higher training and test error than a 20-layer one; the same
phenomenon appears on ImageNet, where a 34-layer \plainnet{} obtains higher
validation error ($28.54\%$ top-1) than its 18-layer counterpart ($27.94\%$
top-1).

The degradation of training accuracy indicates that not all systems are
similarly easy to optimize. Consider a shallower architecture and its deeper
counterpart that adds more layers onto it. There exists a solution by
\emph{construction} to the deeper model: the added layers are identity mappings,
and the other layers are copied from the learned shallower model. The existence
of this constructed solution indicates that a deeper model should produce no
higher training error than its shallower counterpart. Yet experiments show that
our current solvers are unable to find solutions that are comparably good or
better than the constructed solution. Because the \plainnet{}s here use batch
normalization, which ensures forward signals have non-zero variances, and the
backward propagated gradients are verified to exhibit healthy norms, we
\emph{argue} that the difficulty is one of \emph{optimization} and that
vanishing gradients are an \emph{unlikely} cause; we do not, however, formally
prove this, and we conjecture that the deep \plainnet{}s may have exponentially
low convergence rates.

In this paper, we address the \degrad{} by introducing a \emph{deep residual
learning} framework. Instead of hoping each few stacked layers directly fit a
desired underlying mapping, we explicitly let these layers fit a residual
mapping. Formally, denoting the desired underlying mapping as
$\Hmap(\xin)$, we let the stacked nonlinear layers fit another mapping of
$\Fmap(\xin) := \Hmap(\xin) - \xin$. The original mapping is recast into
$\Fmap(\xin) + \xin$. We hypothesize that it is easier to optimize the residual
mapping than to optimize the original, unreferenced mapping. To the extreme, if
an identity mapping were optimal, it would be easier to push the residual to
zero than to fit an identity mapping by a stack of nonlinear layers. The
formulation of $\Fmap(\xin) + \xin$ can be realized by feedforward neural
networks with \emph{shortcut connections} that simply perform identity mapping;
these connections add neither extra parameters nor computational complexity, and
the entire network can still be trained end-to-end by SGD with backpropagation.

We present comprehensive experiments on ImageNet to show the \degrad{} and
evaluate our method. We show that the residual networks are easy to optimize,
whereas the counterpart \plainnet{}s exhibit higher training error when depth
increases; and that the residual networks can easily enjoy accuracy gains from
greatly increased depth, producing results substantially better than previous
networks. Similar phenomena are also shown on the CIFAR-10 set, where successful
models at over 100 layers are presented and a model with over 1000 layers is
explored. On the ImageNet classification dataset, we obtain excellent results by
extremely deep residual nets: a 152-layer residual net is the deepest network
ever presented on ImageNet while still having lower complexity than VGG nets,
and an ensemble has $3.57\%$ top-5 error on the test set, winning the 1st place
in the ILSVRC 2015 classification competition. The extremely deep
representations also have excellent generalization performance on other
recognition tasks: replacing the VGG-16 backbone with a 101-layer residual net
in a baseline detection system yields a $28\%$ relative improvement on the COCO
object detection dataset, indicating that the depth of representations is of
central importance for many visual recognition tasks.
